\chapter{Partitions and a Broker on a Real Network}
The first two chapters encoded records into recoverable single- and multi-segment logs. This chapter gives a topic (a named category of messages) multiple partitions (each partition is its own ordered log), then connects clients to a broker (a process that stores partition data and handles requests) over loopback TCP. Every partition has its own offset sequence: offset 0 in orders-0 and offset 0 in orders-1 identify different records, and they cannot be compared to establish a topic-wide order. Metadata tells clients which partitions exist, their leader broker, and connection endpoint. TCP delivers an ordered byte stream, not Java write boundaries, so an application protocol needs frame lengths. The course protocol is for teaching only and is not compatible with Kafka's wire protocol.

\begin{figure}[ht]
\centering
\begin{tikzpicture}[font=\small,node distance=15mm]
\node[draw,rounded corners,minimum width=30mm,minimum height=10mm] (cli) {client};
\node[draw,rounded corners,minimum width=34mm,minimum height=10mm,right=of cli] (tcp) {TCP frames};
\node[draw,rounded corners,minimum width=38mm,minimum height=10mm,right=of tcp] (broker) {broker handler};
\node[draw,rounded corners,minimum width=42mm,minimum height=10mm,below=of broker] (logs) {partition catalog + logs};
\draw[-{Stealth},thick] (cli) -- node[above]{request} (tcp);
\draw[-{Stealth},thick] (tcp) -- (broker);
\draw[-{Stealth},thick] (broker) -- node[right]{metadata / append / fetch} (logs);
\draw[-{Stealth},thick] (broker.west) |- node[pos=.25,left]{reply} (cli.south);
\end{tikzpicture}
\caption{Request/response path over one connection. Provided code handles message-body encoding and decoding; students implement frames and business dispatch.}
\end{figure}

\lessonsection{9}{Partition Routing}
\goals{Route the same key to the same partition consistently while making null-key rotation observable.}{Understand key bytes, an unsigned CRC32C value, and modulo by a positive partition count.}
\background{A producer chooses a partition before appending so that it knows which ordered log receives a record. The course rule is to compute CRC32C over a non-null key's bytes, interpret the 32-bit result as an unsigned value, and take it modulo the partition count; the same byte string and count therefore choose the same partition. A null key has no bytes to hash, so one Partitioner instance rotates among available partitions. An empty byte array is a present key of length zero and is not the same as null.}
\providededit{The no-argument \method{Partitioner()} constructor is provided, but the round-robin cursor state is not a provided answer; add whatever instance state the implementation needs.}{Implement \method{public int choose(byte[] key,int partitions)} in \pathcode{exercises/src/main/java/io/simplekafka/client/Partitioner.java}.}
\subsubsection*{Prerequisites}
This step depends on \stepref{8}{tail truncation}; the previous chapter established the continuous offsets held by a partition log. The chapter introduction defines topic, partition, and broker: this method chooses an existing partition for a write, but does not create a topic or log.
\subsubsection*{Inputs, outputs, and state}
\begin{itemize}
\item key is null or any non-null byte array, including an empty array; partitions is the number of available partitions and must be >0. Return an int partition id in \pathcode{[0,partitions)}.
\item A non-null key uses \pathcode{unsigned(CRC32C(key)) mod partitions} and does not read or write the null-key cursor; an empty array has CRC32C=0. A null key uses this instance's cursor: normalize \pathcode{cursor mod partitions} to select a partition, then set the cursor to \pathcode{(selected partition+1) mod partitions}; a changed partition count is applied before selection.
\item partitions<=0 reports \method{INVALID_REQUEST} and does not advance the cursor. Different Partitioner instances have independent cursors. This is in-memory state, with no disk or network side effect.
\end{itemize}
\lessoncontract{Validate that the partition count is positive first. Do not use a signed-int remainder when the CRC is meant to be an unsigned 32-bit value. Only a null-key call advances the rotation cursor; keyed and empty-array selections do not. When the partition count changes, each null-key result must remain in range.}
\subsubsection*{Hand calculation: keyed routing and cursor}
Assume key string “123456789” is encoded as ASCII bytes. Its CRC32C is \pathcode{0xe3069283} (unsigned value 3,808,858,755); modulo 7 partitions gives partition 2. An empty array has CRC32C=0, so with 7 partitions it returns partition 0 without changing the null-key cursor. On a fresh instance, null-key calls with partition counts 3, 3, 1, 3, 3, 2, 2 return 0, 1, 0, 0, 1, 0, 1: the count-1 call normalizes the cursor to 0, and later calls use the new partition count.
\expected{\method{Step09Test}: \pathcode{routesByKnownUnsignedCrc32cVectorsAndKeepsEmptyKeyAsAKey} checks CRC32C and empty keys; \pathcode{keyedAndEmptyKeysDoNotAdvanceTheNullKeyCursor} checks cursor preservation; \pathcode{nullKeyCursorsAreIndependentAndSinglePartitionAlwaysSelectsZero} and \pathcode{nullKeyCursorStaysWithinPartitionsWhenTopicSizesChange} check instance isolation and changing partition counts; \pathcode{requiresAPositivePartitionCount} checks invalid counts.}
\pitfalls{Do not use an object's \method{hashCode()} instead of CRC32C over key bytes, and do not take a remainder directly from a possibly negative signed int. Real Kafka has a different key hash and a sticky-batching policy for null keys; the course's per-instance rotation is not Kafka's default behavior.}
\lessonhints{Draw three inputs—null, empty array, and nonempty array—and note which ones change the cursor.}{Interpret CRC as an unsigned number from 0 to $2^{32}-1$, then use a positive partition count for modulo.}{Track the null cursor call by call; normalize when count changes and insert a keyed call to verify it does not advance the cursor.}
\lessoncommand{9}
\lessonquestions{Why should partition hashing use key bytes rather than a RecordData object's Java hashCode?}{How do null and a zero-length key differ in routing and state updates?}
\paragraph*{Self-check explanation}
Hashing uses stable key bytes; Java object hashCode is not the course rule and does not promise the same cross-restart meaning. A null key rotates and updates the instance cursor; an empty array has the actual CRC32C value 0, routes deterministically, and does not update the cursor.

\lessonsection{10}{Topic Metadata}
\goals{Register a fixed set of partitions in the local catalog and return deterministic metadata for client routing.}{Distinguish disk directories from in-memory topic registration.}
\background{A topic is a logical name used by readers; its partition directories contain log files. Metadata is a structured description a broker returns to clients, not a guess based on arbitrary folders on disk. Once the course catalog creates a fixed partition count, a single broker describes each partition with its TopicPartition, current broker as leader, epoch, replicas, ISR, and endpoint. In single-broker mode the epoch is 0, and replicas and ISR contain only that broker.}
\providededit{The \method{PartitionCatalog} constructor fixes root, broker id, and endpoint; \method{TopicPartition}, \method{PartitionMetadata}, log registration, and catalog-root containment are provided.}{Implement \method{public synchronized void createTopic(String topic,int partitionCount)} and \method{public synchronized List<PartitionMetadata> metadata(String topic)} in \pathcode{exercises/src/main/java/io/simplekafka/broker/PartitionCatalog.java}.}
\subsubsection*{Prerequisites}
Complete \stepref{9}{partition routing}; understand that a partition id is scoped within a topic. Also use the \method{PartitionLog} from \stepref{6}{segmented log} as the storage for each partition. Metadata registers/describes partitions; it does not choose a producer's key.
\subsubsection*{Inputs, outputs, and state}
\begin{itemize}
\item topic must contain 1–255 ASCII bytes from \pathcode{[A-Za-z0-9._-]}, and cannot be \pathcode{.} or \pathcode{..}. partitionCount must be positive. Perform these checks before creating topic or partition paths. An invalid name/count reports \method{INVALID_REQUEST} and must not change the directory.
\item Successful \method{createTopic} returns no value, registers partitions 0 through count-1, and opens their logs. Repeating the same name/count is idempotent and does not erase existing data. A different count for the same name reports \method{INVALID_REQUEST} and leaves existing metadata, files, and logs unchanged. Underlying create/open I/O errors report \method{STORAGE_ERROR}. Calling after the catalog is closed throws \method{IllegalStateException}.
\item \method{metadata} returns a list ordered by partition id; in each single-node entry leaderId=current brokerId, epoch=0, replicas=[brokerId], isr=[brokerId], and endpoint is the configured address. An unregistered topic reports \method{UNKNOWN_TOPIC_OR_PARTITION}. A new catalog starts with an empty registry, so explicitly call createTopic to register and reopen the directories; an unexpected numeric directory is not automatically included. Calling after the catalog is closed throws \method{IllegalStateException}.
\end{itemize}
\lessoncontract{Validate the identifier and positive count before creation. Repeating the same definition does not change disk; a different count is rejected without changing the existing directory. Metadata order is stable and does not expose paths outside root. The model permits a topic name of up to 255 bytes; a file system that rejects a 255-byte path component may still fail directory creation, although TopicPartition accepts the identifier itself.}
\subsubsection*{Hand calculation: create, restart, and register}
Assume brokerId=7, a configured endpoint, and an empty root. After createTopic(\pathcode{orders},2), metadata is ordered orders-0 then orders-1, and both entries have leader=7, epoch=0, replicas=isr=[7]. Append records 0 and 1 to orders-0, close the catalog, and create a new one: metadata(\pathcode{orders}) now reports UNKNOWN\_TOPIC\_OR\_PARTITION because the topic registry is not persistent metadata. Calling createTopic(\pathcode{orders},2) again opens the same directory; orders-0 still reads offsets 0 and 1, and a new append returns [2,3). Even an extra on-disk orders/9 directory does not add a third registered entry.
\expected{\method{Step10Test}: \pathcode{createsIndependentFixedTopicsAndReopensOnlyWithExplicitPartitionCounts} checks metadata, persisted logs, and explicit re-registration; \pathcode{rejectsInvalidCountsAndTopicNamesWithoutMutatingTheFilesystemOrEscapingTheRoot} checks validation and unchanged directories; \pathcode{acceptsA255ByteTopicIdentifierWhereTheFilesystemAllowsIt} checks the identifier boundary.}
\pitfalls{Do not concatenate an unvalidated user topic into a path, or enumerate directories to infer a partition count and register manual leftovers. The course catalog does not model Kafka's topic control plane or partition-expansion process.}
\lessonhints{Separate the “topic in the in-memory registry” from “a log directory exists on disk.”}{Write the resulting state for same-name/same-count, same-name/different-count, and unregistered-after-restart calls.}{Compare directory snapshots around invalid creates and restart; after explicit re-registration, check that offsets continue without a gap.}
\lessoncommand{10}
\lessonquestions{Why can't the same topic name with a different partition count be treated as idempotent success?}{After restart, what can an existing numeric directory prove, and what can it not prove?}
\paragraph*{Self-check explanation}
Partition count determines the fixed set of routable partitions; accepting another count silently would make topic metadata and logs disagree. A directory proves some files exist, but does not prove they are registered as the current topic set or reveal the count of partitions that are absent.

\lessonsection{11}{TCP Frame Protocol}
\goals{Split a TCP byte stream into bounded frames and exchange requests or replies one frame at a time.}{Understand short reads, length prefixes, EOF position, and correlation ids.}
\background{TCP provides an ordered byte stream but does not preserve each application's write boundary: one write may arrive through multiple reads, while several writes may be combined in one read. A frame carries a length telling the reader how many bytes make up this protocol unit; the fixed header supplies version, API, request correlation id, and error code, while the payload is the business body. The correlationId is the value used to match a response to its request.}
\providededit{The \method{Frame} value type, API ids, ErrorCode wire ids, \method{Messages} body types, and \method{MessageCodec} body codec are provided; this step does not rewrite business bodies.}{Implement \method{public static void write(OutputStream output,Frame frame) throws IOException} and \method{public static Frame read(InputStream input) throws IOException} in \pathcode{exercises/src/main/java/io/simplekafka/protocol/FrameCodec.java}.}
\subsubsection*{Prerequisites}
Complete \stepref{10}{topic metadata} to know that APIs send requests to a broker, and understand the big-endian fixed-width integers from \stepref{1}{record encoding}. A frame is the course's TCP wrapper, not a disk record or Kafka's official request header.
\subsubsection*{Inputs, outputs, and state}
\begin{itemize}
\item Wire order is big-endian \pathcode{length:int32 | version:int16(1) | api:int16 | correlationId:int32 | error:int16 | payload}. length excludes its own 4 bytes but includes the remaining 10-byte header and payload; legal values are [10,8,388,608], maximum payload is 8,388,598 bytes, and the maximum physical frame is 8,388,612 bytes.
\item \method{write} writes one Frame. Null arguments, unknown APIs, or a total length outside the range report \method{INVALID_REQUEST} before output begins. Success emits exactly one frame. A stream I/O failure propagates the same \method{IOException}; once actual I/O begins, the external stream is not promised to roll back.
\item \method{read} returns at most one frame per call. A null input reports \method{INVALID_REQUEST}. EOF before the first length byte returns null; EOF after only part of the length prefix or frame body reports \method{INVALID_REQUEST}. An out-of-range length, version other than 1, unknown API, or unknown error wire id also reports \method{INVALID_REQUEST}; validate length before reading/allocating its declared size. Input stream I/O failures propagate as the original \method{IOException}.
\end{itemize}
\lessoncontract{A read loop must fill the required bytes rather than assuming one read fills a header or body. Only EOF at a clean frame boundary returns null; a partial frame is not an empty result. Consecutive frames consume their own lengths without swallowing the next frame. For an invalid/oversized prefix, consume no more than its already-read 4-byte length; do not read its suffix or allocate memory based on a malicious size. A known error id in a frame only means its error field is valid; the provided server also rejects a request frame with nonzero error as \method{INVALID_REQUEST} before dispatch.}
\subsubsection*{Decode the golden frame field by field}
The Step11Test request/reply golden is \pathcode{0000000d 0001 0001 00000029 0008 010203}. The first 4 bytes are 13, the length after the prefix: a 10-byte header plus a 3-byte payload. version=\pathcode{0001}=1, API=\pathcode{0001}=METADATA, correlationId=\pathcode{00000029}=41, error=\pathcode{0008}=REQUEST\_TIMEOUT, and payload is \pathcode{010203}. Thus the physical frame is 17 bytes. With no payload, length remains 10 and the full physical frame is 14 bytes.
\expected{\method{Step11Test}: \pathcode{writesAndReadsTheIndependentGoldenFrameAndHandlesEmptyPayload} checks the golden bytes; \pathcode{acceptsMaximumFrameAndRejectsOneByteOverBeforeWritingAnything} checks limits and zero-byte rejection; \pathcode{rejectsEveryIncompletePrefixOfTheGoldenFrame} checks EOF; \pathcode{readsBackToBackFramesFromStreamsDeliveringOneTwoOrThreeBytesAtATime} checks short reads and consecutive frames; \pathcode{rejectsInvalidLengthsWithoutConsumingBytesAfterTheHeader} and \pathcode{rejectsUnknownVersionApiAndErrorIds} check format errors; \pathcode{propagatesTheSameUnderlyingIoExceptionFromReadAndWrite} checks the original I/O exception is propagated.}
\pitfalls{One read is not guaranteed to fill a length or header; allocating from an unvalidated length can amplify memory use for invalid input. The course's version=1, API list, and error ids are teaching wire rules, not a connection to an official Kafka client.}
\lessonhints{Sum the fixed header field widths to get the minimum length 10, and decide whether length includes its own prefix.}{Distinguish “no new frame has started” from “a frame has started but bytes are missing,” and loop to fill short reads.}{Feed the golden one byte at a time to confirm each call consumes only one frame; append a sentinel after an oversized length and verify it is not read.}
\lessoncommand{11}
\lessonquestions{Why should a response echo the correlationId for its request?}{Why is it useful to distinguish clean EOF from EOF halfway through a frame?}
\paragraph*{Self-check explanation}
The correlation id lets a client pair each response with the correct call. Boundary EOF means the stream ended normally; a partial frame is truncated protocol data and cannot be disguised as “no more requests.”
\lessonsection{12}{Single-Broker Requests}
\goals{Connect metadata, produce, and fetch business requests to the local catalog and verify the full path over real loopback TCP.}{Connect frame handling, the provided body codec, and log input/output boundaries.}
\background{A broker is the request-semantics boundary: it dispatches a Request to metadata, produce, or fetch, validates partition, parameters, and epoch, calls the matching backend, and returns a success body or an ErrorBody with an error code. produce appends records to one partition; fetch actively reads records by offset and budget and does not delete the log. A single broker has no follower, so this step's high watermark (HW, the visible committed boundary) equals LEO. A real TCP test exercises client framing, the provided body codec, the server handler, and response framing together.}
\providededit{The \method{BrokerHarness.single}, TCP server/client, \method{MessageCodec} request/reply bodies, request/reply types, log registration, and service lifecycle are provided; \method{SingleBrokerMain} uses loopback and an OS-assigned port.}{Implement \method{public Messages.Reply handle(Messages.Request request)} in \pathcode{exercises/src/main/java/io/simplekafka/broker/BrokerHandler.java}. This step implements METADATA, PRODUCE, and FETCH dispatch only; do not rewrite transport or the body codec.}
\subsubsection*{Prerequisites}
Complete \stepref{11}{TCP frame protocol} and the topic/partition concepts used in \stepref{9}{partition routing} and \stepref{10}{topic metadata}. This step fetches with single-broker HW=LEO; the replication chapter will teach HW again, so do not generalize the single-node rule.
\subsubsection*{Inputs, outputs, and state}
\begin{itemize}
\item Metadata input is \method{MetadataRequest(topic)}; success returns \method{MetadataBody(List<PartitionMetadata>)}. A null or malformed topic returns \method{INVALID_REQUEST}; an unknown but valid topic returns \method{UNKNOWN_TOPIC_OR_PARTITION}. Metadata reads the catalog and does not change logs.
\item Produce input is \method{ProduceRequest(tp,epoch,acks,timeoutMillis,records)}; single-broker mode accepts only epoch=0, acks must be \method{Acks.LEADER} (wire=1) or \method{Acks.ALL} (wire=-1), timeoutMillis must be nonnegative, and records must be nonempty. The single-broker backend ignores the acknowledgment mode, and a nonnegative timeout does not cause a wait; acknowledgment behavior with replicas is taught in \stepref{23}{produce acknowledgments}. Check order: resolve the partition, validate produce parameters, obtain the backend, check epoch/leader, then call backend produce. Unknown partition returns \method{UNKNOWN_TOPIC_OR_PARTITION}; invalid parameters return \method{INVALID_REQUEST}; a mismatched epoch returns \method{FENCED_EPOCH}. Success appends and returns \method{ProduceBody(AppendResult)}. Precondition failures do not append; an I/O failure after writing starts does not guarantee there was no physical side effect.
\item Fetch input is \method{FetchRequest(tp,epoch,offset,maxRecords,maxBytes,token)}; use token=null here, before the later consumer-group fencing step. Check order: resolve the partition, validate positive maxRecords/maxBytes, obtain the backend, check epoch/leader, then fetch from the backend. Single-broker mode accepts epoch=0; an offset outside the retained range \pathcode{[logStartOffset,LEO]} returns \method{OFFSET_OUT_OF_RANGE}, and a request rejected by these checks does not change the log. Success returns \method{FetchBody(records,logStartOffset,logEndOffset,highWatermark,epoch)}; records fit both budgets and are never split. Here HW=LEO and offset=LEO is an empty successful read. Fetch does not change the log.
\item The handler returns a \method{Reply}: success has error=\pathcode{NONE} and its matching success body; failure carries a course error code and safe \method{ErrorBody}. A null or unsupported request reports \method{INVALID_REQUEST}. The handler maps a \method{CourseException} to a Reply preserving its code; parameter \method{IllegalArgumentException}/\method{NullPointerException} maps to \method{INVALID_REQUEST}; unexpected runtime failures map to \method{STORAGE_ERROR} without exposing their cause. \method{ExerciseNotImplementedException} passes through so the course runner can locate unfinished work. Provided transport handles invalid frames, nonzero request-frame errors, body format/trailing bytes, and reply-frame encoding; the handler must not reimplement these parsing tasks.
\end{itemize}
\lessoncontract{Unknown partitions, stale epochs, and out-of-range offsets retain their own error codes; a request rejected before backend access must not append. Single-broker metadata has epoch=0, and a successful fetch does not pass LEO=HW. An error affects only its request: do not send Java stack traces or private exception details over the wire, and keep handling later valid TCP requests. The provided \method{MessageCodec} caps fetch maxBytes at 8,388,566 disk bytes to fit a single response frame; an oversized complete record is not split, and remaining records are fetched in a later request starting at the next offset.}
\subsubsection*{Hand calculation: produce, fetch, and page a large reply}
Assume orders-0 exists, epoch=0, LEO=0, and produce uses acks=\method{Acks.LEADER} and timeoutMillis=0. Append three records in order: key \pathcode{k} (1 byte)/value \pathcode{v0} (2 bytes), key \pathcode{k}/value \pathcode{v1}, and null key/value \pathcode{v2}. Ordinary record totals are 35, 35, and 34 bytes. Success returns [0,3) and LEO becomes 3. A fetch at offset=0 with maxRecords=2 and maxBytes=4096 returns offsets 0 and 1; FetchBody has start=0, LEO=3, HW=3, epoch=0. Fetch at offset=3 returns empty. If maxBytes is 69, one byte below the total for two 35-byte records, only offset 0 is returned.
\par Now assume 9 records, each with a null key and a 1,048,548-byte value; each ordinary disk record is 1,048,580 bytes. A request for maxBytes=10,000,000 is capped by the provided codec to 8,388,566: 7 records use 7,340,060 bytes and fit, so the first reply contains offsets 0–6; 8 records use 8,388,640 bytes and exceed the cap, so offset 7 is not included. A later request starting at offset 7 returns 7 and 8. Records are never split to fill a budget.
\expected{\method{Step12Test}: \pathcode{servesPreciseMetadataRecordsAndFailuresWithoutMutatingTheLog} checks metadata, produce/fetch, and rejected requests leaving the log unchanged; \pathcode{boundsLargeFetchRepliesAndPreservesRecordContinuation} and \pathcode{returnsAnErrorReplyWhenAHandlerProducesAnOversizedFetchBody} check large-response budgets; \pathcode{sanitizesHandlerRuntimeExceptionsAndKeepsBothConnectionsUsable} checks error sanitization and continued service; \pathcode{isolatesExistingTopicPartitionsAndServesFetchesFromNonzeroLogStart} checks partition/start isolation; \pathcode{rawTcpRequestsAndMalformedBodiesStayIsolatedToTheirConnections} checks failure isolation over real TCP.}
\pitfalls{Calling only the handler directly misses TCP boundaries, correlation ids, and reply-codec errors; swallowing an exception as an empty success makes a caller believe a write/read succeeded. One business error must not stop the broker. A successful produce and the client receiving its reply are separate network events: a reply may be unreachable after a network failure.}
\lessonhints{List input type, prechecks, and success body for METADATA, PRODUCE, and FETCH as three separate paths.}{Keep the responsibility and error code of each frame/MessageCodec, handler-validation, and backend failure layer distinct.}{Use loopback TCP for metadata, produce, and fetch, then send a known failing request and fetch again; compare the log and replies.}
\lessoncommand{12}
\lessonquestions{Why must an error Reply not send a Java stack trace or internal exception message to a client?}{Why does single-broker HW=LEO not define the rule for replica mode?}
\paragraph*{Self-check explanation}
Internal exceptions may disclose implementation details and are not a stable protocol contract; return only a course error code and short safe message. A single broker has no replication progress to distinguish; with replicas, LEO is only a local tail and HW separately defines the visible range.
